\section{Circuit Complexity --- Strict Fan Out}
%%% generated by the blueprint pipeline from TCSlib.Complexity.CircuitComplexity.StrictFanOut
\subsection*{Circuit Complexity --- Strict Fan Out: Introduction}

Declarations of the Lean module \texttt{TCSlib.Complexity.CircuitComplexity.StrictFanOut}.

\subsection*{Circuit Complexity --- Strict Fan Out: Results}

\begin{theorem}[Value of a negation gate]
\lean{BoolCircuit.vertexValue_not_gate}
\label{BoolCircuit.vertexValue_not_gate}
\leanok
\uses{BoolCircuit.DAGGate, BoolCircuit.DAGGate.eval, BoolCircuit.runWith_getD_gate, BoolCircuit.runWith_getD_take, BoolCircuit.vertexValue}
\difficulty{3}
Let $gs$ be a list of gates, let $i$ and $w$ be natural numbers, and let $x$ assign
Boolean values to the inputs. Suppose $i<|gs|$, the gate at position $i$ is a negation
gate reading vertex $w$, and $w<n+i$, where $n$ is the number of inputs. Then the value
at vertex $n+i$ is the negation of the value at vertex $w$.
\end{theorem}

\begin{theorem}[Counting a shifted range]
\lean{BoolCircuit.countP_range_add}
\label{BoolCircuit.countP_range_add}
\leanok
\difficulty{4}
For natural numbers $c,u,k$, the number of indices $j$ with $0 \le j < k$ and $u=c+j$ is
$1$ if $c \le u < c+k$, and is $0$ otherwise.
\end{theorem}

\begin{theorem}[Lookup in an Appended List]
\lean{BoolCircuit.getElem_append_of_eq}
\label{BoolCircuit.getElem_append_of_eq}
\leanok
\difficulty{3}
Let $l_1$ and $l_2$ be lists of elements of a type $\alpha$, and let $i,j$ be natural
numbers. If $i=|l_1|+j$, $j<|l_2|$, and $i<|l_1\mathbin{++}l_2|$, then the entry at
index $i$ of the concatenated list is the entry at index $j$ of $l_2$.
\end{theorem}

\begin{definition}[One-step gate transformation]
\lean{BoolCircuit.nnStep}
\label{BoolCircuit.nnStep}
\leanok
\uses{BoolCircuit.DAGGate}
For $n\in\mathbb N$, a state $(G,t,N)$ consisting of a finite list $G$ of gates, a map
$t\colon\mathbb N\to\mathbb N$, and a counter $N\in\mathbb N$, and a gate $g$ with
argument list of length $k$, define the next state as follows. Append to $G$ a gate of
the same kind as $g$ whose arguments are the images under $t$ of the arguments of $g$,
then one negation gate reading each such image, then $k$ negation gates reading, in
order, the vertices $n+|G|+1$ through $n+|G|+k$. The new map sends $N$ to $n+|G|$; for
any other vertex appearing among the arguments of $g$, it sends that vertex to
$n+|G|+1+k$ plus the index of its first occurrence in the argument list; on all
remaining vertices it agrees with $t$. The counter increases to $N+1$.
\end{definition}

\begin{definition}[Iterated fan-out reduction]
\lean{BoolCircuit.nnGates}
\label{BoolCircuit.nnGates}
\leanok
\uses{BoolCircuit.DAGGate, BoolCircuit.nnStep}
Given a natural number $n$ and a list of gates, start with the empty gate list, the
identity vertex map, and counter $n$, then apply the one-step fan-out reduction
successively to each gate in the given list. The resulting gate list, vertex map, and
counter form the output.
\end{definition}

\begin{definition}[Invariant for fan-out reduction]
\lean{BoolCircuit.NNInv}
\label{BoolCircuit.NNInv}
\leanok
\uses{BoolCircuit.DAGGate, BoolCircuit.DAGGate.Strict, BoolCircuit.DAGGate.eval, BoolCircuit.GatesAcyclic, BoolCircuit.fanoutIn, BoolCircuit.runWith, BoolCircuit.vertexValue}
For a natural number $n$, an original gate list, and a state consisting of a gate list,
a vertex map, and a counter, the invariant requires that the counter equal $n$ plus the
length of the original list; the new gate list be acyclic; and every original vertex
index map below $n$ plus the original list length to a vertex below $n$ plus the new
gate-list length, with distinct original indices mapping to distinct vertices. It also
requires that, for every assignment of Boolean values to the $n$ input vertices,
evaluation at each mapped original vertex agrees with the output of the original circuit
on that assignment; every vertex have fan-out at most two; each mapped original vertex
have fan-out zero; every new gate have distinct arguments and have exactly one argument
if it is a negation gate and exactly two otherwise; and the new list have length at most
five times the original list length.
\end{definition}

\begin{theorem}[Invariant for the Empty Gate List]
\lean{BoolCircuit.nnInv_nil}
\label{BoolCircuit.nnInv_nil}
\leanok
\uses{BoolCircuit.GatesAcyclic.nil, BoolCircuit.NNInv, BoolCircuit.fanoutIn}
\difficulty{3}
For every natural number $n$, the empty gate list, identity tap map, and counter $n$
satisfy the construction invariant: the counter equals $n$, the constructed gate list is
acyclic, taps on vertices below $n$ are below $n$ and injective, and every tapped vertex
has the same value as the corresponding vertex of the empty computation for every
Boolean input. Each vertex has fanout at most two, each tapped vertex has fanout zero,
all constructed gates satisfy the required distinct-input and arity conditions, and the
number of constructed gates is at most zero.
\end{theorem}

\begin{theorem}[Preservation of the circuit invariant under one step]
\lean{BoolCircuit.nnInv_step}
\label{BoolCircuit.nnInv_step}
\leanok
\uses{BoolCircuit.DAGGate, BoolCircuit.DAGGate.Strict, BoolCircuit.DAGGate.Strict.faninTwo, BoolCircuit.DAGGate.eval, BoolCircuit.DAGGate.eval_remap, BoolCircuit.GatesAcyclic, BoolCircuit.NNInv, BoolCircuit.countP_range_add, BoolCircuit.fanoutIn, BoolCircuit.fanoutIn_append, BoolCircuit.fanoutIn_eq_zero_of_le, BoolCircuit.gatesAcyclic_append, BoolCircuit.getElem_append_of_eq, BoolCircuit.nnStep, BoolCircuit.runWith, BoolCircuit.runWith_append, BoolCircuit.runWith_getD_last, BoolCircuit.runWith_getD_of_lt, BoolCircuit.vertexValue_append, BoolCircuit.vertexValue_last, BoolCircuit.vertexValue_not_gate}
\difficulty{7}
Let $n$ be a natural number, let $G$ be a finite list of Boolean gates, let $g$ be a
Boolean gate, and let $s$ be a state. If $s$ satisfies the invariant for $G$, the list
obtained by appending $g$ to $G$ is acyclic relative to $n$, and $g$ has distinct
arguments with exactly one argument when it is a NOT gate and exactly two arguments when
it is an AND or OR gate, then the state produced by applying one construction step to
$s$ using $g$ satisfies the invariant for the list obtained by appending $g$ to $G$.
\end{theorem}

\begin{theorem}[Invariants of the monotone gate construction]
\lean{BoolCircuit.nnGates_inv}
\label{BoolCircuit.nnGates_inv}
\leanok
\uses{BoolCircuit.DAGGate, BoolCircuit.DAGGate.Strict, BoolCircuit.GatesAcyclic, BoolCircuit.NNInv, BoolCircuit.gatesAcyclic_append, BoolCircuit.nnGates, BoolCircuit.nnInv_nil, BoolCircuit.nnInv_step}
\difficulty{4}
For every natural number $n$ and every finite list $G$ of gates, if the gates in $G$ are
acyclic relative to $n$ and each gate has distinct inputs and exactly one input for a
NOT gate or two for an AND or OR gate, then the construction obtained by processing $G$
has the following properties. Its counter equals $n+|G|$; its gates are acyclic; its tap
map sends every index below $n+|G|$ to an index below $n$ plus the number of constructed
gates, and is injective on those indices. For every Boolean assignment to the $n$
inputs, each such tapped vertex has the same value as the corresponding vertex in the
original computation; every vertex is used by at most two constructed gates, and tapped
vertices are unused. Every constructed gate has distinct inputs and the prescribed arity
for its kind, and the number of constructed gates is at most $5|G|$.
\end{theorem}

\begin{definition}[Fan-out-two circuit conversion]
\lean{BoolCircuit.DAGCircuit.fanoutTwoNN}
\label{BoolCircuit.DAGCircuit.fanoutTwoNN}
\leanok
\uses{BoolCircuit.DAGCircuit, BoolCircuit.DAGGate.Strict, BoolCircuit.nnGates, BoolCircuit.nnGates_inv}
Let $n$ be a natural number, and let $C$ be a directed acyclic circuit with $n$ input
vertices. Suppose every gate has distinct input vertices, with one input for a negation
gate and two inputs for a conjunction or disjunction gate. The fan-out reduction
construction applied to $C$ defines another directed acyclic circuit with the same $n$
input vertices.
\end{definition}

\begin{theorem}[Evaluation preserved by bounded-fanout conversion]
\lean{BoolCircuit.DAGCircuit.fanoutTwoNN_eval}
\label{BoolCircuit.DAGCircuit.fanoutTwoNN_eval}
\leanok
\uses{BoolCircuit.DAGCircuit.eval, BoolCircuit.DAGCircuit.fanoutTwoNN, BoolCircuit.nnGates_inv}
\difficulty{1}
Let $n$ be a nonnegative integer, and let $C$ be a directed acyclic Boolean circuit on
$n$ inputs whose gates each have distinct arguments and exactly one argument for a NOT
gate or two for an AND or OR gate. For every input assignment, the circuit obtained by
the bounded-fanout conversion computes the same Boolean value as $C$.
\end{theorem}

\begin{theorem}[Strict gates after fanout reduction]
\lean{BoolCircuit.DAGCircuit.fanoutTwoNN_strict}
\label{BoolCircuit.DAGCircuit.fanoutTwoNN_strict}
\leanok
\uses{BoolCircuit.DAGCircuit.fanoutTwoNN, BoolCircuit.DAGGate.Strict, BoolCircuit.nnGates_inv}
\difficulty{1}
Let $C$ be a finite directed acyclic Boolean circuit, and suppose each gate reads
distinct input vertices, with exactly one input for a NOT gate and exactly two for an
AND or OR gate. Every gate in the circuit obtained by the fanout reduction construction
also has distinct input vertices, with exactly one input for a NOT gate and exactly two
for an AND or OR gate.
\end{theorem}

\begin{theorem}[Fanout bound after normalization]
\lean{BoolCircuit.DAGCircuit.fanoutTwoNN_hasFanoutTwo}
\label{BoolCircuit.DAGCircuit.fanoutTwoNN_hasFanoutTwo}
\leanok
\uses{BoolCircuit.DAGCircuit.HasFanoutTwo, BoolCircuit.DAGCircuit.fanoutTwoNN, BoolCircuit.nnGates_inv}
\difficulty{1}
Let $C$ be a finite directed acyclic Boolean circuit, and suppose each gate reads
distinct input vertices, with exactly one input for a NOT gate and exactly two for an
AND or OR gate. Then in the circuit obtained by the fanout reduction construction, every
vertex is read by at most two gates.
\end{theorem}

\begin{theorem}[Output has zero fanout after conversion]
\lean{BoolCircuit.DAGCircuit.fanoutTwoNN_fanout_output}
\label{BoolCircuit.DAGCircuit.fanoutTwoNN_fanout_output}
\leanok
\uses{BoolCircuit.DAGCircuit.fanout, BoolCircuit.DAGCircuit.fanoutTwoNN, BoolCircuit.nnGates_inv}
\difficulty{1}
Let $n$ be a nonnegative integer, and let $C$ be a directed acyclic Boolean circuit on
$n$ inputs whose gates each have distinct arguments and exactly one argument for a NOT
gate or two for an AND or OR gate. In the circuit obtained by the bounded-fanout
conversion, the output vertex has fanout zero.
\end{theorem}

\begin{theorem}[Size bound after fanout reduction]
\lean{BoolCircuit.DAGCircuit.fanoutTwoNN_size_le}
\label{BoolCircuit.DAGCircuit.fanoutTwoNN_size_le}
\leanok
\uses{BoolCircuit.DAGCircuit.fanoutTwoNN, BoolCircuit.DAGCircuit.size, BoolCircuit.nnGates, BoolCircuit.nnGates_inv}
\difficulty{3}
Let $C$ be a finite directed acyclic Boolean circuit, and suppose each gate reads
distinct input vertices, with exactly one input for a NOT gate and exactly two for an
AND or OR gate. The size of the circuit obtained by the fanout reduction construction is
at most $n+5m$, where $n$ is the number of input vertices and $m$ is the number of gates
of $C$.
\end{theorem}

\begin{theorem}[Bounded-fanout circuit conversion]
\lean{BoolCircuit.DAGCircuit.exists_fanoutTwo_notNot}
\label{BoolCircuit.DAGCircuit.exists_fanoutTwo_notNot}
\leanok
\uses{BoolCircuit.DAGCircuit, BoolCircuit.DAGCircuit.HasFanoutTwo, BoolCircuit.DAGCircuit.eval, BoolCircuit.DAGCircuit.fanoutTwoNN, BoolCircuit.DAGCircuit.fanoutTwoNN_eval, BoolCircuit.DAGCircuit.fanoutTwoNN_hasFanoutTwo, BoolCircuit.DAGCircuit.fanoutTwoNN_size_le, BoolCircuit.DAGCircuit.fanoutTwoNN_strict, BoolCircuit.DAGCircuit.size, BoolCircuit.DAGGate.Strict}
\difficulty{3}
Let $n$ be a nonnegative integer, and let $C$ be a directed acyclic Boolean circuit on
$n$ inputs whose gates each have distinct arguments and exactly one argument for a NOT
gate or two for an AND or OR gate. There is a circuit $C'$ on the same inputs with these
same gate conditions, with fanout at most $2$ at every vertex, and computing the same
Boolean function as $C$ on every input assignment, such that $C'$ has size at most $5$
times the size of $C$.
\end{theorem}

\begin{theorem}[Bounded-fanout conversion for strict circuits]
\lean{BoolCircuit.DAGCircuit.exists_isStrict_fanoutTwo_of_strict}
\label{BoolCircuit.DAGCircuit.exists_isStrict_fanoutTwo_of_strict}
\leanok
\uses{BoolCircuit.DAGCircuit, BoolCircuit.DAGCircuit.HasFanoutTwo, BoolCircuit.DAGCircuit.IsStrict, BoolCircuit.DAGCircuit.eval, BoolCircuit.DAGCircuit.fanoutTwoNN, BoolCircuit.DAGCircuit.fanoutTwoNN_eval, BoolCircuit.DAGCircuit.fanoutTwoNN_fanout_output, BoolCircuit.DAGCircuit.fanoutTwoNN_hasFanoutTwo, BoolCircuit.DAGCircuit.fanoutTwoNN_size_le, BoolCircuit.DAGCircuit.fanoutTwoNN_strict, BoolCircuit.DAGCircuit.seal, BoolCircuit.DAGCircuit.seal_eval, BoolCircuit.DAGCircuit.seal_hasFanoutTwo, BoolCircuit.DAGCircuit.seal_isStrict, BoolCircuit.DAGCircuit.seal_size_le, BoolCircuit.DAGCircuit.size, BoolCircuit.DAGGate.Strict}
\difficulty{4}
Let $n$ be a nonnegative integer, and let $C$ be a directed acyclic Boolean circuit on
$n$ inputs in which every gate has distinct arguments and exactly one argument for a NOT
gate or two for an AND or OR gate. There is a circuit $C'$ on the same inputs that
computes the same Boolean function, in which every gate has those same distinctness and
arity properties, every vertex other than the output has positive fanout, and every
vertex has fanout at most $2$; moreover, $C'$ has size at most $20$ times the size of
$C$.
\end{theorem}

\begin{theorem}[Bounded-fan-out equivalent circuit]
\lean{BoolCircuit.DAGCircuit.IsStrict.exists_fanoutTwo}
\label{BoolCircuit.DAGCircuit.IsStrict.exists_fanoutTwo}
\leanok
\uses{BoolCircuit.DAGCircuit, BoolCircuit.DAGCircuit.HasFanoutTwo, BoolCircuit.DAGCircuit.IsStrict, BoolCircuit.DAGCircuit.eval, BoolCircuit.DAGCircuit.exists_isStrict_fanoutTwo_of_strict, BoolCircuit.DAGCircuit.size}
\difficulty{1}
Let $C$ be a finite Boolean circuit in which every gate has distinct arguments and has
exactly one argument if it is a NOT gate, and exactly two otherwise; suppose every
vertex other than the output has positive fan-out. Then there exists another circuit
$C'$ with the same number of inputs such that every gate of $C'$ satisfies those same
argument conditions, every vertex has fan-out at most two, $C'$ computes the same
Boolean function as $C$ on every input, and the size of $C'$ is at most $20$ times the
size of $C$.
\end{theorem}

\begin{theorem}[Strict bounded-fanout circuit conversion]
\lean{BoolCircuit.DAGCircuit.exists_isStrict_fanoutTwo}
\label{BoolCircuit.DAGCircuit.exists_isStrict_fanoutTwo}
\leanok
\uses{BoolCircuit.DAGCircuit, BoolCircuit.DAGCircuit.HasFanoutTwo, BoolCircuit.DAGCircuit.IsFaninTwo, BoolCircuit.DAGCircuit.IsStrict, BoolCircuit.DAGCircuit.deconst, BoolCircuit.DAGCircuit.deconst_eval, BoolCircuit.DAGCircuit.deconst_size, BoolCircuit.DAGCircuit.deconst_strict, BoolCircuit.DAGCircuit.eval, BoolCircuit.DAGCircuit.exists_isStrict_fanoutTwo_of_strict, BoolCircuit.DAGCircuit.size}
\difficulty{4}
Let $n$ be a positive integer, and let $C$ be a directed acyclic Boolean circuit on $n$
inputs in which every gate has distinct arguments and at most two arguments, and every
NOT gate has exactly one argument. There is a circuit $C'$ on the same inputs that
computes the same Boolean function, in which every gate has distinct arguments and
exactly one argument for a NOT gate or two for an AND or OR gate, every vertex other
than the output has positive fanout, and every vertex has fanout at most $2$; moreover,
$C'$ has size at most $20$ times the size of $C$ plus $60$.
\end{theorem}
